\chapter{Prezentacja warstwy użytkowej projektu}
\label{cha:prezentacja}

W tym rozdziale przedstawiono wizualną stronę aplikacji za pomocą zrzutów ekranu, ilustrujących
kluczowe widoki oraz interakcje użytkownika. Celem jest zaprezentowanie funkcjonalności systemu
z perspektywy użytkownika końcowego - zarówno nauczyciela, jak i ucznia. Kolejne podrozdziały
odpowiadają typowemu scenariuszowi pracy: nauczyciel loguje się, tworzy i udostępnia test, uczeń
rozwiązuje przypisany mu test, a następnie nauczyciel ocenia odpowiedzi i analizuje wyniki.

% =====================================================================
\section{Ekran logowania i rejestracji}
\label{sec:logowanie}

Pierwszym ekranem, z którym styka się użytkownik, jest ekran logowania, umożliwiający dostęp do
systemu zarejestrowanym użytkownikom lub utworzenie nowego konta (ucznia lub nauczyciela). Interfejs
przedstawiono na rysunku~\ref{rys:logowanie}. Składa się on z:

\begin{itemize}
  \item pola \emph{Login}, służącego do wprowadzenia nazwy użytkownika,
  \item pola \emph{Hasło}, zapewniającego maskowanie wpisywanego hasła,
  \item przycisku \emph{Zaloguj się}, który weryfikuje podane dane z bazą danych i w przypadku ich
        poprawności przenosi użytkownika do odpowiedniego panelu (nauczyciela lub ucznia),
  \item zakładki \emph{Rejestracja}, umożliwiającej utworzenie nowego konta po podaniu imienia,
        nazwiska, loginu i hasła oraz wybraniu roli (ucznia lub nauczyciela),
  \item komunikatów zwrotnych, informujących o błędnych danych logowania lub o problemach
        z walidacją formularza rejestracji.
\end{itemize}

\figscreen{ekran_logowanie.png}{0.75}{Ekran logowania i rejestracji.}{rys:logowanie}

% =====================================================================
\section{Pulpit nauczyciela}
\label{sec:pulpit}

Po zalogowaniu na konto nauczyciela wyświetlany jest pulpit (rysunek~\ref{rys:pulpit}),
stanowiący przegląd najważniejszych informacji. Zawiera on:

\begin{itemize}
  \item kafelki podsumowujące liczbę testów, uczniów i grup oraz średni wynik,
  \item listę testów nauczyciela z szybkim przejściem do ich statystyk,
  \item panel ostatnich wyników uczniów,
  \item boczny panel nawigacji umożliwiający przejście do pozostałych modułów: testów, oceniania,
        grup i statystyk.
\end{itemize}

\figscreen{ekran_pulpit.png}{0.95}{Pulpit nauczyciela.}{rys:pulpit}

% =====================================================================
\section{Lista testów i edytor testu}
\label{sec:edytor}

Widok \emph{Testy} prezentuje wszystkie testy nauczyciela w formie tabeli z informacjami o nazwie
i statusie widoczności (aktywny/nieaktywny), przedmiocie, liczbie pytań, czasie, przypisanych grupach
i liczbie podejść. Status można przełączać jednym kliknięciem, a przy każdym teście dostępne są
akcje: statystyki, edycja, duplikowanie oraz usunięcie. Z poziomu tego widoku nauczyciel tworzy nowy
test lub otwiera istniejący w edytorze (rysunek~\ref{rys:edytor}).

Edytor testu pozwala określić ustawienia testu (nazwę, przedmiot wpisywany dowolnie lub wybierany
z listy, limit czasu, próg zaliczenia, liczbę podejść, widoczność dla uczniów oraz przypisane grupy)
i zbudować listę pytań. Dla każdego pytania można wybrać typ (jednokrotny wybór, wielokrotny wybór
lub pytanie otwarte), wpisać treść, ustalić liczbę punktów oraz - w pytaniach zamkniętych -
warianty odpowiedzi i poprawne odpowiedzi. Edytor waliduje kompletność testu przed zapisem (m.in.
obecność poprawnej odpowiedzi i przypisanej grupy). Jeżeli test ma już rozwiązania, edycja pytań
jest blokowana, o czym informuje stosowny komunikat.

\figscreen{ekran_edytor.png}{0.95}{Lista testów nauczyciela ze statusem widoczności i akcjami.}{rys:listaTestow}

\figscreen{ekran_testy.png}{0.95}{Edytor testu wraz z ustawieniami i listą pytań.}{rys:edytor}

% =====================================================================
\section{Ocenianie pytań otwartych}
\label{sec:ocenianie}

Moduł \emph{Ocenianie} (rysunek~\ref{rys:ocenianie}) gromadzi podejścia uczniów zawierające pytania
otwarte oczekujące na ręczną ocenę. Dla każdego zgłoszenia nauczyciel widzi:

\begin{itemize}
  \item dane ucznia oraz nazwę testu i datę oddania,
  \item treść pytania otwartego oraz odpowiedź wpisaną przez ucznia,
  \item podpowiedź w postaci odpowiedzi wzorcowej,
  \item pole do wpisania przyznanych punktów w dozwolonym zakresie (z walidacją zakresu).
\end{itemize}

Po przyznaniu punktów wszystkim pytaniom otwartym i zapisaniu oceny wynik ucznia jest automatycznie
przeliczany, a zgłoszenie znika z kolejki.

\figscreen{ekran_ocenianie.png}{0.9}{Ocenianie pytań otwartych.}{rys:ocenianie}

% =====================================================================
\section{Zarządzanie grupami}
\label{sec:grupy}

Widok \emph{Grupy} (rysunek~\ref{rys:grupy}) prezentuje grupy uczniów w formie kart. Każda karta
zawiera nazwę i opis grupy, listę przypisanych uczniów oraz listę testów udostępnionych grupie.
Przycisk zarządzania otwiera okno, w którym nauczyciel zaznacza uczniów należących do grupy.
Dodatkowo, z poziomu tego widoku można utworzyć nową grupę.

\figscreen{ekran_grupy.png}{0.95}{Widok zarządzania grupami uczniów.}{rys:grupy}

% =====================================================================
\section{Statystyki testu i eksport wyników}
\label{sec:statystyki}

Moduł \emph{Statystyki} (rysunek~\ref{rys:statystyki}) prezentuje analizę wyników wybranego testu.
Zawiera:

\begin{itemize}
  \item wskaźnik zdawalności oraz liczbę uczniów, łączną liczbę podejść i średni wynik,
  \item informację o najtrudniejszym pytaniu,
  \item wykres poprawności w rozbiciu na poszczególne pytania,
  \item tabelę z wszystkimi podejściami uczniów (najlepsze podejście oznaczone gwiazdką) wraz ze
        statusem (zaliczony, niezaliczony, oczekujący na ocenę, próba oszustwa) oraz możliwością
        ręcznej edycji oceny każdego podejścia,
  \item przycisk eksportu wyników do pliku CSV.
\end{itemize}

Wskaźniki zbiorcze (zdawalność, średni wynik) liczone są z najlepszego podejścia każdego ucznia.

\figscreen{ekran_statystyki.png}{0.95}{Statystyki wyników testu.}{rys:statystyki}

% =====================================================================
\section{Lista testów ucznia}
\label{sec:testyUcznia}

Po zalogowaniu na konto ucznia wyświetlana jest lista testów przypisanych do jego grup
(rysunek~\ref{rys:testyUcznia}), z podziałem na testy do zrobienia oraz ukończone. Każdy test
prezentowany jest w formie karty z informacją o przedmiocie, liczbie pytań, czasie, liczbie
wykorzystanych podejść oraz - dla testów rozwiązanych - z osiągniętym wynikiem i możliwością jego
podglądu lub ponownego podejścia (jeśli limit nie został wyczerpany, a test jest nadal aktywny).
Uczeń widzi wyłącznie testy aktywne oraz te nieaktywne, które już rozwiązał (zachowując dostęp do
swojego wyniku); po dezaktywacji testu niewykorzystane podejścia nie są już dostępne.

\figscreen{ekran_testy_ucznia.png}{0.95}{Lista testów przypisanych uczniowi.}{rys:testyUcznia}

% =====================================================================
\section{Rozwiązywanie testu}
\label{sec:rozwiazywanie}

Ekran rozwiązywania testu (rysunek~\ref{rys:rozwiazywanie}) prezentuje pojedyncze pytanie wraz
z wariantami odpowiedzi lub polem na odpowiedź otwartą. W górnej części znajduje się licznik czasu,
a u dołu - nawigacja umożliwiająca przechodzenie między pytaniami oraz pasek postępu. Po upływie
czasu test jest automatycznie kończony. Przy oddawaniu testu z pominiętymi pytaniami wyświetlane
jest ostrzeżenie. Opuszczenie okna testu jest traktowane jako próba oszustwa.

\figscreen{ekran_rozwiazywanie.png}{0.95}{Ekran rozwiązywania testu.}{rys:rozwiazywanie}

% =====================================================================
\section{Wynik testu}
\label{sec:wynik}

Po zakończeniu testu uczeń widzi ekran wyniku (rysunek~\ref{rys:wynik}), prezentujący osiągnięty
procent, liczbę zdobytych punktów oraz status (zaliczony, niezaliczony lub oczekujący na ocenę
pytań otwartych). Poniżej znajduje się szczegółowy przegląd odpowiedzi - dla każdego pytania
wskazane są odpowiedzi poprawne oraz wybrane przez ucznia.

\figscreen{wynik_testu.png}{0.95}{Ekran wyniku testu wraz z przeglądem odpowiedzi.}{rys:wynik}
